\begin{ouverture}[statistiques]{19}{Données et enquêtes}{Partie 4 · Données, hasard et instructions}
{Question · Recueil · Tableau · Lecture critique}
\objectif{Planifier une enquête et recueillir des données.}
\objectif{Consigner des observations et construire un tableau simple.}
\objectif{Lire des représentations, filtrer des données et argumenter un choix.}
\prerequis{Tableaux, diagrammes en barres, diagrammes circulaires et courbes du cours moyen.}
\end{ouverture}
\pagestyle{otMain}\essentiel\fontsize{11}{15}\selectfont
\phantomsection\label{cours-19-1}
\essSection{1}{Préparer une enquête}
\begin{savoirfaire}[Les étapes d’une enquête]1. Formuler une question précise. 2. Choisir qui interroger ou quoi observer. 3. Recueillir les données : questionnaire, comptage ou mesures. 4. Les organiser dans un tableau. 5. Les représenter et les interpréter. 6. Répondre à la question, avec prudence.\end{savoirfaire}
\begin{definition}Le caractère étudié peut être une qualité (moyen de transport, couleur préférée) ou une quantité (temps de trajet, taille d’une plante). L’effectif d’une valeur est le nombre de fois où elle apparaît.\end{definition}
\essSection{2}{Mesurer et consigner des observations}
\begin{exemple}24 élèves ont indiqué leur moyen de transport. On compte les réponses et on construit le tableau des effectifs ; le total contrôle le comptage.\begin{center}\renewcommand{\arraystretch}{1.35}\begin{tabular}{c|ccccc}\toprule
Moyen de transport & à pied & vélo & bus & voiture & Total\\\midrule
Effectif & 9 & 4 & 8 & 3 & 24\\
\bottomrule\end{tabular}\end{center}\end{exemple}
\begin{savoirfaire}[Consigner des mesures]Pour des mesures, on note la date, l’unité et l’instrument utilisé, puis on range les valeurs dans un tableau, une colonne par observation.\end{savoirfaire}
\clearpage
\phantomsection\label{cours-19-2}
\essSection{3}{Lire tableaux et représentations}
\label{fig-cours-19-2-cours-graphiques}\figureGeo{cours-graphiques}

\begin{propriete}Un diagramme en barres compare des effectifs ; un diagramme circulaire montre la part de chaque catégorie dans le tout ; une courbe montre une évolution, souvent au cours du temps.\end{propriete}
\essSection{4}{Filtrer des données et argumenter un choix}
\begin{savoirfaire}[Filtrer]Pour faire un choix, on garde seulement les lignes d’un tableau qui vérifient un premier critère (jour, prix…), puis on applique le critère suivant aux lignes restantes.\end{savoirfaire}
\begin{attention}Lire avec esprit critique : d’où viennent les données ? Combien de personnes ont été interrogées ? L’axe vertical commence-t-il à zéro ? Une conclusion ne vaut que pour les données recueillies.\end{attention}

\clearpage\rubrique{Découverte}
\begin{blocLivre}{ex-19-1}\begin{exercice}[1]{Le bon outil}
Pour chaque question d’enquête, dire si l’on utilise plutôt un questionnaire, un comptage ou des mesures.\begin{questions}
\item Quel est le sport préféré des élèves de la classe ?
\item Combien de véhicules passent devant le collège entre 8 h et 8 h 30 ?
\item Quelle est la taille des élèves de la classe ?
\end{questions}
\end{exercice}\end{blocLivre}
\begin{blocLivre}{ex-19-2}\begin{exercice}[1]{Construire un tableau}
Vingt élèves ont répondu à la question « Quel animal as-tu à la maison ? » : chat, chien, chat, poisson, aucun, chat, chien, aucun, chat, lapin, chien, chat, aucun, chat, poisson, chien, chat, aucun, lapin, chat. Construire le tableau des effectifs.
\end{exercice}\end{blocLivre}
\begin{blocLivre}{ex-19-3}\begin{exercice}[1]{Lire un diagramme en barres}
Le diagramme indique le nombre de livres empruntés au CDI chaque jour. Quel jour a-t-on emprunté le plus de livres ? Combien de livres en tout dans la semaine ? Quel écart entre le jour le plus fort et le plus faible ?
\label{fig-ex-19-3}\figureGeo{don-3}
\end{exercice}\end{blocLivre}
\begin{blocLivre}{ex-19-4}\begin{exercice}[1]{Lire un diagramme circulaire}
Le diagramme représente les 4 heures d’écran d’un adolescent un dimanche. Combien de temps a-t-il passé sur chaque activité ?
\label{fig-ex-19-4}\figureGeo{don-4}
\end{exercice}\end{blocLivre}
\begin{blocLivre}{ex-19-5}\begin{exercice}[1]{Lire une courbe}
La courbe montre la température relevée dans une cour de récréation au cours d’une journée de printemps. Quelle est la température la plus élevée, et à quelle heure ? De combien la température a-t-elle augmenté entre 6 h et ce moment ?
\label{fig-ex-19-5}\figureGeo{don-5}
\end{exercice}\end{blocLivre}

\rubrique{Entraînement}
\begin{blocLivre}{ex-19-6}\begin{exercice}[2]{Relevé de mesures}
Un plant de haricot mesure 2 cm la semaine 1, puis 5 cm, $9$ cm, $14{,}5$ cm et 17 cm les semaines suivantes. Présenter ces mesures dans un tableau. Entre quelles semaines a-t-il le plus grandi ?
\end{exercice}\end{blocLivre}
\begin{blocLivre}{ex-19-7}\begin{exercice}[2]{Tableau à double entrée}
Le tableau donne le moyen de transport des élèves de deux classes. Compléter les totaux. Combien d’élèves viennent à pied ? Combien d’élèves en tout ?\begin{center}\renewcommand{\arraystretch}{1.35}\begin{tabular}{c|cccc}\toprule
Classe & à pied & vélo & bus & Total\\\midrule
6e A & 10 & 4 & 12 & …\\
6e B & 8 & 6 & 11 & …\\
Total & … & … & … & …\\
\bottomrule\end{tabular}\end{center}
\end{exercice}\end{blocLivre}
\begin{blocLivre}{ex-19-8}\begin{exercice}[2]{Filtrer pour choisir}
Sami cherche un atelier le mercredi, à 10 € au plus, avec encore des places. Quels ateliers peut-il choisir ?\begin{center}\renewcommand{\arraystretch}{1.35}\begin{tabular}{c|ccc}\toprule
Atelier & Jour & Prix (€) & Places restantes\\\midrule
Échecs & mercredi & 8 & 3\\
Théâtre & mercredi & 15 & 5\\
Robotique & samedi & 9 & 2\\
Chorale & mercredi & 0 & 0\\
Escalade & mercredi & 9 & 4\\
\bottomrule\end{tabular}\end{center}
\end{exercice}\end{blocLivre}
\begin{blocLivre}{ex-19-9}\begin{exercice}[2]{Un graphique trompeur}
Un magasin affiche ce diagramme : « Nos ventes explosent ! » Les ventes sont passées de 102 à 108 articles. Pourquoi le dessin est-il trompeur ?
\label{fig-ex-19-9}\figureGeo{don-9}
\end{exercice}\end{blocLivre}
\begin{blocLivre}{ex-19-10}\begin{exercice}[2]{Construire un diagramme}
Construire le diagramme en barres du tableau suivant, avec une graduation de 2 en 2.\begin{center}\renewcommand{\arraystretch}{1.35}\begin{tabular}{c|cccc}\toprule
Sport & foot & danse & judo & natation\\\midrule
Effectif & 12 & 9 & 5 & 4\\
\bottomrule\end{tabular}\end{center}
\label{fig-ex-19-10}\figureGeo{don-10}
\end{exercice}\end{blocLivre}

\rubrique{Problèmes}
\begin{blocLivre}{ex-19-11}\begin{exercice}[3]{Le gaspillage à la cantine}
Les délégués pèsent le pain jeté à la cantine : $3{,}2$ kg lundi, $2{,}8$ kg mardi, $3{,}5$ kg jeudi et $2{,}5$ kg vendredi. Quelle masse de pain est jetée dans la semaine ? Le collège veut réduire ce gaspillage de 25~\% : quelle masse viser ?
\end{exercice}\end{blocLivre}
\begin{blocLivre}{ex-19-12}\begin{exercice}[3]{Les insectes du jardin}
Une classe a compté les insectes observés dans le jardin du collège. Quelles espèces ont été vues au moins 3 jours sur 5 ? Quelle espèce a été comptée le plus souvent en tout ?\begin{center}\renewcommand{\arraystretch}{1.35}\begin{tabular}{c|ccccc}\toprule
Espèce & lun & mar & mer & jeu & ven\\\midrule
Abeille & 4 & 0 & 3 & 5 & 2\\
Coccinelle & 0 & 1 & 0 & 2 & 0\\
Papillon & 2 & 3 & 1 & 0 & 4\\
Fourmi & 10 & 12 & 8 & 15 & 9\\
Sauterelle & 0 & 0 & 1 & 0 & 0\\
\bottomrule\end{tabular}\end{center}
\end{exercice}\end{blocLivre}
\begin{blocLivre}{ex-19-13}\begin{exercice}[3]{Le climat d’une ville}
Le graphique donne, pour une ville, la température moyenne (courbe) et les précipitations (barres) de chaque mois. Quel est le mois le plus chaud ? Le plus sec ? Le plus pluvieux ? Que remarque-t-on en été ?
\label{fig-ex-19-13}\figureGeo{don-13}
\end{exercice}\end{blocLivre}
\begin{blocLivre}{ex-19-14}\begin{exercice}[3]{Planifier une enquête}
La classe veut répondre à la question : « Combien de temps les élèves mettent-ils pour venir au collège ? » Décrire les étapes de l’enquête et proposer un tableau vide pour les résultats.
\end{exercice}\end{blocLivre}
\begin{blocLivre}{ex-19-15}\begin{exercice}[3]{Choisir un séjour}
Une famille cherche un séjour au bord de la mer, de 5 jours au moins et à 400 € au plus. Lequel peut-elle choisir ?\begin{center}\renewcommand{\arraystretch}{1.35}\begin{tabular}{c|ccc}\toprule
Séjour & Prix (€) & Durée (jours) & Lieu\\\midrule
Alpes & 380 & 5 & montagne\\
Bretagne & 420 & 6 & mer\\
Vendée & 350 & 5 & mer\\
Corse & 520 & 7 & mer\\
Landes & 390 & 4 & mer\\
\bottomrule\end{tabular}\end{center}
\end{exercice}\end{blocLivre}

\clearpage
